\subsection{Phonon--phonon scattering self-energy}
\label{sub:phph_se}
% ----------------------------------------------------------------------

The anharmonic self-energy has been reported in slightly different forms in the
literature: the prefactor reported for the same diagram in~\citet{wang_nonequilibrium_2007,wang_quantum_2008},
\citet{mingo_anharmonic_2006}, and \citet{luisier_atomistic_2012} differ by
factors of $\pi/9$ and $4$, as analysed by~\citet{guo_quantum_2020}
who also provided the diagrammatic derivation.
We give a derivation below.

Let $\hat{V}_3(\tau) = \tfrac{1}{3!}\sum_{\mu\nu\lambda}
\Phi_{\mu\nu\lambda}\,\tilde{u}_\mu(\tau)\,\tilde{u}_\nu(\tau)\,\tilde{u}_\lambda(\tau)$
be the cubic interaction in the interaction picture (With respect to the harmonic operators). For the interaction
we compute the S-matrix~\cite{rammer_quantum_1986,haug_quantum_2008}
\begin{equation}\label{eq:S_matrix}
  S_{\mathcal{C}} = T_{\mathcal{C}}\exp\!\left[-\frac{i}{\hbar}\oint_{\mathcal{C}}d\tau\,\hat{V}_3(\tau)\right].
\end{equation}
Wick's theorem (valid for the contour-ordered Green's function~\cite{wang_quantum_2008}) decomposes contour-ordered
products into pair contractions. The full Green's function is
\begin{equation}\label{eq:G_full}
  G_{\mu\nu}(\tau,\tau')
  = -\frac{i}{\hbar}\,\frac{
       \langle T_{\mathcal{C}}\,
       S_{\mathcal{C}}\,
       \tilde{u}_\mu(\tau)\,
       \tilde{u}_\nu(\tau')\rangle_0}
      {\langle S_{\mathcal{C}}\rangle_0},
\end{equation}
and the linked-cluster theorem cancels the denominator against the
disconnected diagrams in the numerator, leaving only connected diagrams
in the self-energy.

% ----------------------------------------------------------------------
%  (a) Sunset diagram
% ----------------------------------------------------------------------
\begin{figure}[ht]
\centering
\begin{tikzpicture}[every node/.style={font=\small}]
  \coordinate (v1) at (0,0);  \coordinate (v2) at (4,0);
  \draw[phonon] (-1.8,0) -- (v1);
  \draw[phonon] (v2)      -- (5.8,0);
  \node[above] at (-1.0, 0.05) {$\mu,\tau$};
  \node[above] at ( 4.95,0.05) {$\mu',\tau'$};
  \draw[phonon] (v1) to[bend left=55]  (v2);
  \draw[phonon] (v1) to[bend right=55] (v2);
  \node at (2, 1.65) {$G_{\nu_1\nu_4}$};
  \node at (2,-1.65) {$G_{\nu_2\nu_3}$};
  \node[fc3 vertex] at (v1) {};
  \node[fc3 vertex] at (v2) {};
  \node[below=18pt] at (v1) {$\Phi_{\mu\nu_1\nu_2}$};
  \node[below=18pt] at (v2) {$\Phi_{\mu'\nu_3\nu_4}$};
\end{tikzpicture}
\caption{Sunset (two-vertex bubble) Feynman diagram for the leading
$(\Phi^{(3)})^2$ phonon--phonon self-energy of
\cref{eq:sigma_sunset_contour}. Wavy lines denote phonon propagators
$G$, filled circles denote cubic FC3 vertices. The two equivalent
internal lines produce the symmetry factor $\tfrac{1}{2}$.}
\label{fig:sunset_diagram}
\end{figure}


To order $(\Phi^{(3)})^2$, the proper self-energy is given by the
``sunset'' diagram (\cref{fig:sunset_diagram}), in which two cubic
vertices are connected by two internal phonon propagators. The
value of the diagram, including all combinatorial factors,
is~\cite{guo_quantum_2020,wang_nonequilibrium_2014}
\begin{equation}\label{eq:sigma_sunset_contour}
  \Sigma^{(2)}_{\mu\mu'}(\tau,\tau')
  = \frac{i\hbar}{2}\sum_{\nu_1\nu_2\nu_3\nu_4}
    \Phi_{\mu\nu_1\nu_2}\,
    G_{\nu_1\nu_4}(\tau,\tau')\,
    G_{\nu_2\nu_3}(\tau,\tau')\,
    \Phi_{\mu'\nu_3\nu_4}.
\end{equation}

\begin{figure}[ht]
\centering
\begin{tikzpicture}[every node/.style={font=\small}]
  \coordinate (v) at (0,0);
  \draw[phonon] (-2.5,0) -- (v);
  \draw[phonon] (v)      -- (2.5,0);
  \node[above] at (-2.0, 0.05) {$\mu,\tau$};
  \node[above] at ( 2.0, 0.05) {$\mu',\tau'$};
  \draw[phonon] (v) to[bend left=80, looseness=2.0] (0,2.4);
  \draw[phonon] (0,2.4) to[bend left=80, looseness=2.0] (v);
  \node at (0,2.95) {$G^{<}_{\nu\lambda}$};
  \node[fc4 vertex] at (v) {};
  \node[below=12pt] at (v) {$\Phi^{(4)}_{\mu\mu'\nu\lambda}$};
\end{tikzpicture}
\caption{Hartree--Fock (one-loop) tadpole self-energy from a single
quartic FC4 vertex. The filled square denotes
$\Phi^{(4)}_{\mu\mu'\nu\lambda}$. The two internal lines contract into
the closed loop $G^{<}_{\nu\lambda}(\tau,\tau)$. The diagram is real
and frequency-independent and acts as a temperature-dependent
renormalisation $D\mapsto D + \Sigma^{(4),\mathrm{HF}}$.}
\label{fig:fc4_tadpole}
\end{figure}


A second diagram, lower order in derivatives but one loop, survives:
the FC4 Hartree--Fock ``loop'' tadpole
% REVIEW: tadpole prefactor corrected so the result is real. With the equal-time
% relation <u_nu u_lambda> = i*hbar*G^<_{nu lambda}(tau,tau), the Hartree tadpole
% (1/2)Phi^(4)<uu> equals (i*hbar/2)Phi^(4)G^<; the previous -hbar/2*G^< was
% imaginary. Confirm the overall sign against static_self_energy.py.
$\Sigma^{(4),\mathrm{HF}}_{\mu\mu'} = \tfrac{1}{2}
\sum_{\nu\lambda} \Phi^{(4)}_{\mu\mu'\nu\lambda}\,\langle u_\nu u_\lambda\rangle
= \tfrac{i\hbar}{2}
\sum_{\nu\lambda} \Phi^{(4)}_{\mu\mu'\nu\lambda}\,
G^{<}_{\nu\lambda}(\tau,\tau)$ from a single quartic vertex (\cref{fig:fc4_tadpole}).
The loop self-energy is real and frequency-independent.
It therefore produces no phonon linewidth, but shifts the phonon
dispersion, $D \mapsto D + \Sigma^{(4),\mathrm{HF}}$, providing the
leading temperature-dependent mode renormalisation of self-consistent
phonon theory~\cite{tadano_first-principles_2018,wang_nonequilibrium_2014}. In a
transport-focused SCBA targeting the imaginary part of $\Sigma_s$, the
loop can be absorbed into a renormalised harmonic spectrum
computed independently (e.g.\ via SCP or
SSCHA~\cite{paulatto_first-principles_2015,tadano_first-principles_2018}) and is not
re-included diagram-by-diagram in the iteration.

The factor $1/2$ above corresponds to the fully-symmetric FC3
$\Phi_{\mu\nu\lambda}$ defined by
$\Phi_{\mu\nu\lambda} = \partial^3 V/(\partial u_\mu \partial u_\nu \partial u_\lambda)$,
divided by $3!$ in \cref{eq:H_phonon}. ~\citet{mingo_anharmonic_2006} obtains an extra factor of $\pi/9$
which~\citet{guo_quantum_2020} attribute to a different convention in
the symmetrization of the cubic vertex.

Applying the Langreth analytic continuation rules~\cite{devreese_linear_1976,haug_quantum_2008}
to the contour product of two $G$'s in
\cref{eq:sigma_sunset_contour} yields real-time lesser/greater
components,
\begin{equation}\label{eq:sigma_guo}
  \Sigma^{\gtrless}_{\mu\mu'}(\omega)
  = \frac{i\hbar}{2}
    \sum_{\nu_1\nu_2\nu_3\nu_4}
    \Phi_{\mu\nu_1\nu_2}
    \left[
      \int\frac{d\omega'}{2\pi}\,
      G^{\gtrless}_{\nu_1\nu_4}(\omega')\,
      G^{\gtrless}_{\nu_2\nu_3}(\omega-\omega')
    \right]
    \Phi_{\mu'\nu_3\nu_4},
\end{equation}
which is the form derived in detail in~\cite{wang_nonequilibrium_2014,guo_quantum_2020}.
The convolution enforces the
energy-conservation $\delta(\omega - \omega' - \omega'')$ implicit in
the contour integral. Substituting equilibrium FDT
\cref{eq:phonon_fdt} reproduces Fermi's-golden-rule expressions for
three-phonon scattering rates: in mode basis,
% REVIEW: absorption-channel Bose factor corrected to 2[n_2 - n_1]. On the
% delta(omega_s - omega_{s1} + omega_{s2}) channel, energy conservation forces
% omega_{s1} > omega_{s2} (so n_1 < n_2); the Maradudin-Fein / phono3py linewidth
% and the bubble combination (1+n_1)n_2 - n_1(1+n_2) = n_2 - n_1 are positive.
\begin{align}\label{eq:fgr}
  \mathrm{Im}\,\Sigma^R_s(\omega_s,\qp)
  &\;=\;
   -\,\frac{\pi\hbar}{8 \omega_s(\qp)}
   \sum_{s_1 s_2}\sum_{\qp{}'}
   |V_3(s\qp; s_1\qp{}'; s_2,\qp{-}\qp{}')|^2 \notag\\
  &\quad\times
   \Bigl\{
     [1 + n_1 + n_2]\,\delta(\omega_s - \omega_{s_1} - \omega_{s_2})
     + 2[n_2 - n_1]\,\delta(\omega_s - \omega_{s_1} + \omega_{s_2})
   \Bigr\},
\end{align}
with $n_i = n_{\mathrm{BE}}(\omega_{s_i},T)$, exhibiting the
splitting/recombination ($1\to 2$) channel and the absorption ($2\to 1$)
channel. This expression coincides with the Boltzmann result of~\citet{klemens_anharmonic_1966}
in the relaxation-time approximation.

\subsubsection{Retarded self-energy}
\label{subsub:KK}

By construction $\Sigma^R$ is causal in time, so its frequency
components satisfy a Kramers--Kronig relation. Starting from
% REVIEW: sign fixed. The Keldysh identities are Sigma^t + Sigma^{tbar} =
% Sigma^< + Sigma^>, Sigma^t - Sigma^{tbar} = Sigma^R + Sigma^A; the retarded
% combination is (1/2)(Sigma^t - Sigma^{tbar}) + (1/2)(Sigma^> - Sigma^<).
% (With the previous '+' the RHS collapsed to Sigma^>.) The final eq:sigma_R_full
% below is unaffected.
$\Sigma^R = \tfrac{1}{2}(\Sigma^t - \Sigma^{\bar{t}}) +
\tfrac{1}{2}(\Sigma^> - \Sigma^<)$ and using the spectral identity for
$\Sigma^{t,\bar{t}}$, one obtains~\cite{guo_quantum_2020,wang_nonequilibrium_2014}
\begin{equation}\label{eq:sigma_R_full}
  \Sigma^R(\omega)
  = \tfrac{1}{2}\bigl[\Sigma^>(\omega) - \Sigma^<(\omega)\bigr]
    + i\,\mathcal{P}\!\!\int_{-\infty}^{\infty}\!\frac{d\omega'}{2\pi}\,
      \frac{\Sigma^>(\omega') - \Sigma^<(\omega')}{\omega - \omega'}.
\end{equation}
The antihermitian (dissipative) first term carries inverse-lifetime
information,
% REVIEW: units fixed. With the omega^2 spectral variable (eq:system), Sigma^R has
% units (rad/s)^2 and the linewidth (HWHM in omega) is -Im Sigma^R/(2 omega_s), so
% the inverse lifetime (FWHM rate) is tau^{-1} = -Im Sigma^R/omega_s -- no factor
% of hbar. (was 2/(hbar omega_s), dimensionally not a rate.) Consistent with the
% companion shift Delta omega = Re Sigma^R/(2 omega_s) below.
\begin{equation}\label{eq:lifetime}
  \tau_s^{-1}(\qp)
  \equiv -\frac{1}{\omega_s(\qp)}\,
     \mathrm{Im}\,\langle s|\Sigma^R(\omega_s,\qp)|s\rangle,
\end{equation}
while the dispersive second term encodes the anharmonic frequency
renormalization $\Delta\omega_s \approx
\mathrm{Re}\,\Sigma^R/(2\omega_s)$ (mode hardening or softening).
Numerical evaluation of the principal value is expensive, so in practice
the dispersive part is often dropped:
\begin{equation}\label{eq:sigma_r}
  \Sigma^R(\omega)
  \approx \tfrac{1}{2}\bigl[\Sigma^>(\omega) - \Sigma^<(\omega)\bigr].
\end{equation}

\paragraph{Half rule versus the causal transform.}
Dropping the dispersive term is harmless at finite $\eta$ and on short devices,
where $\mathrm{Re}\,\Sigma^R$ is a small renormalisation and the lead broadening
keeps the propagator well conditioned. At $\eta\to0$ on longer devices it is not.
The half rule \cref{eq:sigma_r} keeps only the anti-Hermitian part, so its
$\Sigma^R$ is no longer the causal Kramers--Kronig partner of the bubble
$\Sigma^{\gtrless}$; the propagator $G^<=G^R\Sigma^<G^A$ built from it is then not
the causal resolvent, and although the \emph{lead-summed} balance of
\cref{eq:sumrule} still holds (the Keldysh route makes $D(\omega)\equiv0$), the
\emph{per-interface} currents \cref{eq:I_local} become non-uniform and can turn
negative, so the $\eta=0$ fixed point is lost. Restoring the full transform
\cref{eq:sigma_R_full} --- evaluated by FFT through the Hilbert transform of the
$\Sigma^{\gtrless}$ data --- recovers a causal $\Sigma^R$, physical per-interface
currents, and a stable conserving fixed point. The half rule is therefore the
finite-$\eta$ default, while the FFT/Kramers--Kronig form is mandatory for the
$\eta\to0$ limit of \cref{sub:eta_zero}.

\paragraph{Storage sign convention.}
The production solver stores \emph{occupation-positive} propagators,
$-iG^{\lessgtr}(\omega)\succeq0$ for $\omega>0$, the same sign as the lead
injection $\Sigma_\ell^{\lessgtr}=+i\,n_\ell(\pm1)\Gamma_\ell$
(\cref{sec:solver_units}). The bubble \cref{eq:sigma_guo} is quadratic in $G$
and therefore returns a \emph{textbook-signed} self-energy,
$-i\Sigma^{\lessgtr}\preceq0$ --- the opposite sign to what the Keldysh feedback
\cref{eq:keldysh_eq} expects in this convention. The implementation accordingly
negates $\Sigma^{\lessgtr}$ on output, while the dissipative retarded part
\cref{eq:sigma_r} is formed from the \emph{raw} (un-negated) values so that the
damping $\Gamma_\Sigma=i(\Sigma^R-\Sigma^A)\succeq0$ keeps its sign. With the
wrong sign the feedback is anti-dissipative and the SCBA diverges at any
coupling; the single-shot positivity gate \texttt{sigma\_convention}
(\cref{app:conservation}) verifies $-i\Sigma^{\lessgtr}\succeq0$ and catches
exactly this error, which the sign-blind energy balance \cref{eq:bubble_balance}
cannot.

\subsubsection{Detailed balance}

The $\Sigma^{\gtrless}$ of \cref{eq:sigma_guo} satisfy detailed balance
at equilibrium,
\begin{equation}\label{eq:detailed_balance}
  \Sigma^>(\omega) = e^{\beta\hbar\omega}\,\Sigma^<(\omega),
\end{equation}
which together with the Keldysh equation \cref{eq:keldysh_eq}
guarantees that $G^{\lessgtr}$ also satisfy detailed balance, recovering
\cref{eq:phonon_fdt} when $\Sigma_L,\Sigma_R$ correspond to a single
temperature. The presence of two leads at different temperatures breaks
detailed balance and drives a steady-state heat current.

\subsubsection{Higher-order corrections}
\label{subsub:higher_order}
\begin{figure}[ht]
\centering
\begin{tikzpicture}[every node/.style={font=\small}]
  \coordinate (v1) at (0,0);     \coordinate (v2) at (5,0);
  \coordinate (vt) at (2.5, 1.4);  \coordinate (vb) at (2.5,-1.4);
  \draw[phonon] (-1.8,0) -- (v1);
  \draw[phonon] (v2)      -- (6.8,0);
  \node[above] at (-1.0, 0.05) {$\mu,\tau$};
  \node[above] at ( 5.95,0.05) {$\mu',\tau'$};
  \draw[phonon] (v1) -- (vt);   \draw[phonon] (v1) -- (vb);
  \draw[phonon] (vt) -- (v2);   \draw[phonon] (vb) -- (v2);
  \draw[phonon] (vt) -- (vb);
  \foreach \pt in {v1,v2,vt,vb}{\node[fc3 vertex] at (\pt) {};}
\end{tikzpicture}
\caption{Two-loop $\Phi^{(3)}$ vertex correction to the phonon
self-energy ($\mathcal{O}((\Phi^{(3)})^4)$). Four cubic vertices form a
1PI ``diamond'' subgraph.}
\label{fig:vertex_correction}
\end{figure}

\begin{figure}[ht]
\centering
\begin{subfigure}[b]{0.95\linewidth}
\centering
\begin{tikzpicture}[every node/.style={font=\small}]
  \coordinate (v1) at (0  ,0);  \coordinate (v2) at (1.6,0);
  \coordinate (v3) at (3.2,0);  \coordinate (v4) at (4.8,0);
  \draw[phonon] (-1.8,0) -- (v1);
  \draw[phonon] (v4)      -- (6.6,0);
  \node[above] at (-1.0,0.05) {$\mu,\tau$};
  \node[above] at ( 5.7,0.05) {$\mu',\tau'$};
  \draw[phonon] (v1) -- (v2);
  \draw[phonon] (v2) -- (v3);
  \draw[phonon] (v3) -- (v4);
  \draw[phonon] (v1) to[bend left=70, looseness=1.1] (v3);
  \draw[phonon] (v2) to[bend left=70, looseness=1.1] (v4);
  \foreach \pt in {v1,v2,v3,v4}{\node[fc3 vertex] at (\pt) {};}
\end{tikzpicture}
\caption{Crossed two-loop self-energy graph at order
$(\Phi^{(3)})^4$.}
\label{fig:crossed_graph}
\end{subfigure}

\vspace{1em}
\begin{subfigure}[b]{0.95\linewidth}
\centering
\begin{tikzpicture}[every node/.style={font=\small}]
  \coordinate (v1) at (0  ,0);   \coordinate (v2) at (1.8,0);
  \coordinate (v3) at (3.4,0);   \coordinate (v4) at (5.2,0);
  \draw[phonon] (-1.8,0) -- (v1);
  \draw[phonon] (v4)      -- (7.0,0);
  \node[above] at (-1.0,0.05) {$\mu,\tau$};
  \node[above] at ( 6.1,0.05) {$\mu',\tau'$};
  \draw[phonon] (v1) to[bend left=80]  (v2);
  \draw[phonon] (v1) to[bend right=80] (v2);
  \draw[phonon] (v3) to[bend left=80]  (v4);
  \draw[phonon] (v3) to[bend right=80] (v4);
  \draw[phonon] (v2) -- (v3);
  \foreach \pt in {v1,v2,v3,v4}{\node[fc3 vertex] at (\pt) {};}
\end{tikzpicture}
\caption{Ladder topology at order $(\Phi^{(3)})^4$:
two sunset bubbles connected by a single line.}
\label{fig:ladder_graph}
\end{subfigure}

\caption{Non-trivial two-loop self-energy topologies at order
$(\Phi^{(3)})^4$ that lie beyond bubble-SCBA.}
\label{fig:two_loop_topologies}
\end{figure}

\Cref{eq:sigma_guo} captures the leading three-phonon process. Higher
orders include~\cite{wang_nonequilibrium_2014,guo_quantum_2020}:
(i) the four-phonon tadpole from $\Phi^{(4)}$,
% REVIEW: prefactor matches the corrected real form of the tadpole defined above.
$\Sigma^{(4),\mathrm{HF}}_{\mu\mu'} = \tfrac{i\hbar}{2}\sum_{\nu\lambda}
\Phi^{(4)}_{\mu\mu'\nu\lambda}\,G^<_{\nu\lambda}(\tau,\tau)$,
which is real, frequency-independent, and renormalizes
$D \mapsto D + \Sigma^{(4),\mathrm{HF}}$, contributing the
temperature-dependent thermal expansion and frequency
shifts~\cite{wang_nonequilibrium_2014} (\cref{fig:fc4_tadpole}),
(ii) the two-loop $\Phi^{(3)}$ vertex correction, relevant at high
$T$ where $\hbar\omega_{\max}/k_B T \lesssim 1$ and three-phonon
processes saturate (\cref{fig:vertex_correction}) and
(iii) crossed and ladder graphs (\cref{fig:two_loop_topologies}) that are needed
to maintain $\Phi$-derivability.

The $\Phi$-derivability framework of Baym and Kadanoff~\cite{baym_conservation_1961, baym_self-consistent_1962}
states that $\Sigma$ must be the functional derivative of a generating
functional $\Phi[G]$ for the conservation laws (energy, particle number)
to be satisfied identically. The sunset diagram of
\cref{eq:sigma_sunset_contour} is $\Phi$-derivable when the same
$G$ enters both internal lines, but only asymptotically as the
SCBA iteration converges~\cite{lee_anharmonic_2018}. This is e.g. addressed by
the LOA-based scheme of~\citet{lee_anharmonic_2018}
which is exactly conserving at every iteration.

\subsubsection{Transverse periodic case}

For a 3D system with periodicity transverse to transport, the
self-energy acquires a $\qp$ argument and an internal sum over
transverse momenta~\cite{guo_quantum_2020}:
\begin{align}\label{eq:sigma_phph}
  \sse{\gtrless}_{x x'}(\omega,\qp)
  = \frac{i\hbar}{2}
    &\sum_{y_1 y_2 y_3 y_4}
    \frac{1}{N}\sum_{\qp{}'}
    \tilde{\Phi}_{x y_1 y_2}\!\bigl(\qp{}',\qp - \qp{}'\bigr)\,
    \tilde{\Phi}_{x' y_3 y_4}\!\bigl(\qp{}'-\qp,-\qp{}'\bigr)
    \notag\\
  &\times
    \int\!\dw\;
    G^{\gtrless}_{y_1 y_4}(\omega',\qp{}')\,
    G^{\gtrless}_{y_2 y_3}(\omega-\omega',\qp - \qp{}'),
\end{align}
where $x,y$ denote transport-direction DOF indices, i.e.\ $x=(l_x,\kappa,\alpha)$
with $l_x$ enumerating transport-direction slabs/supercells (here $(\kappa,\alpha)$
is the intra-slab DOF, distinct from the global index $\mu=(l,\kappa,\alpha)$ of
\cref{eq:H_phonon}). The
partially Fourier-transformed FC3 entering this expression is
\begin{align}\label{eq:fc3_fourier}
  \tilde{\Phi}_{x y_1 y_2}(\qp,\qp{}')
  = \tilde{\Phi}^{l m_1 m_2}_{\mu \nu_1 \nu_2}(\qp,\qp{}')
  = \sum_{\substack{n_1,\,n_2\\ n_{ix}=m_{ix}}}
    \Phi^{l n_1 n_2}_{\mu \nu_1 \nu_2}\;
    e^{-i\qp\cdot (n_1 - m_1)}\;
    e^{-i\qp{}'\cdot (n_2 - m_2)},
\end{align}
i.e.\ the sum runs over transverse supercells at fixed transport-direction
cell. Unlike the harmonic FC2 object which depends on a single
transverse wavevector, the cubic tensor depends on two: the
$\qp{}'$ sum couples all transverse wavevectors, so evaluating
$\Sigma_s(\qp)$ requires Green's functions at every $\qp{}'$, destroying
the embarrassing parallelism of the ballistic problem. The arithmetic
complexity scales as $\mathcal{O}(N_q^2 N_\omega^2 N_{\mathrm{DOF}}^4)$
(the $N_{\mathrm{DOF}}^4$, rather than the naive $N_{\mathrm{DOF}}^6$ of the
four-index ring, already assumes the FC3 block-locality of
\cref{ssec:phph_sparsity})
before any low-rank or stochastic acceleration. This is the dominant
cost of the anharmonic simulation, as reported in~\cite{luisier_atomistic_2012,guo_quantum_2020, lee_anharmonic_2018}.

% ----------------------------------------------------------------------
\subsection{Observables}
\label{sub:negf_obs}
% ----------------------------------------------------------------------

We collect the observables computable from $G^{R,A,\lessgtr}$ in both
the ballistic and anharmonic regimes.

\subsubsection{Spectral function and local density of states}

The spectral function and LDOS are
\begin{equation}\label{eq:dos}
  A(\omega^2) = i[G^R - G^A],
  \qquad
  \rho_\mu(\omega^2) = \frac{1}{\pi}\,A_{\mu\mu}(\omega^2),
\end{equation}
related to the standard $\omega$-resolved density of states by
$\rho(\omega) = 2\omega\,\rho(\omega^2)$. In the transverse periodic
case $\rho_\mu(\omega) = (2\omega/N_q)\sum_{\qp} A_{\mu\mu}(\omega^2,\qp)/\pi$.

\subsubsection{Ballistic transmission and conductance}

In the ballistic limit ($\Sigma_s = 0$) the transmission function is
given by the Caroli formula~\cite{caroli_direct_1971,mingo_calculation_2003,wang_nonequilibrium_2007},
\begin{equation}\label{eq:caroli}
  \mathcal{T}(\omega^2,\qp)
  = \mathrm{Tr}\bigl[\Gamma_L\,G^R\,\Gamma_R\,G^A\bigr],
\end{equation}
with the transverse-averaged transmission
\begin{equation}
  \bar{\mathcal{T}}(\omega^2)
  = \frac{1}{N_q}\sum_{\qp} \mathcal{T}(\omega^2,\qp).
\end{equation}
The thermal conductance per unit area is the Landauer integral
\begin{equation}\label{eq:G_thermal}
  G(T)
  = \frac{1}{A_c}\int_0^\infty\!\frac{d\omega}{2\pi}\;
    \bar{\mathcal{T}}(\omega)\,\hbar\omega\,
    \frac{\partial n_\mathrm{BE}}{\partial T},
\end{equation}
with $A_c$ the cross-sectional area and
\begin{equation}
  \frac{\partial n_\mathrm{BE}}{\partial T}
  = \frac{\hbar\omega}{k_B T^2}\,
    \frac{e^{\hbar\omega/k_B T}}
         {(e^{\hbar\omega/k_B T}-1)^2}.
\end{equation}
The integrand vanishes at both endpoints, so a trapezoidal rule on the
sampled frequency grid suffices.

\subsubsection{Mode-resolved transmission}

Diagonalizing $\Gamma_\ell$ at fixed $(\omega,\qp)$ yields incoming and
outgoing mode bases of the leads. Writing $\Gamma_L = U_L \Lambda_L
U_L^\dagger$ and similarly for $R$, the transmission matrix element from
incoming mode $m\in L$ to outgoing mode $m'\in R$ is
\begin{equation}\label{eq:mode_T}
  t_{m'm}(\omega,\qp)
  = \sqrt{\Lambda_R^{m'}}\,
    \bigl[U_R^\dagger G^R U_L\bigr]_{m'm}\,
    \sqrt{\Lambda_L^{m}},
\end{equation}
with $\sum_{m'm}|t_{m'm}|^2 = \mathcal{T}$. This is the phonon analogue
of the Fisher--Lee relation~\cite{fisher_relation_1981}.
For phonons the mode-matching formulation~\cite{ong_efficient_2015}
provides an alternative route to the same
quantity through the eigenproblem of the bulk transfer matrix.

\subsubsection{Meir--Wingreen heat current}

In the anharmonic regime the Caroli formula no longer holds: there is
no $\mathcal{T}(\omega)$ that captures both elastic and inelastic
processes. The generalization is the phonon Meir--Wingreen
formula~\cite{meir_landauer_1992,wang_nonequilibrium_2007,wang_quantum_2008},
\begin{equation}\label{eq:meir_wingreen}
  I_L
  = \int_0^\infty\!\frac{d\omega}{4\pi}\,\hbar\omega\,
    \mathrm{Tr}\bigl[
      \Sigma_L^<(\omega)\,G^>(\omega)
    - \Sigma_L^>(\omega)\,G^<(\omega)
    \bigr].
\end{equation}
\Cref{eq:meir_wingreen} is exact for any
many-body $\Sigma_s$, provided $G^{\lessgtr}$ are obtained from the full
Keldysh equation \cref{eq:keldysh_eq}. In the ballistic limit, the
Keldysh equation gives $G^{<} = G^R(\Sigma_L^< + \Sigma_R^<)G^A$, and
substitution into \cref{eq:meir_wingreen} reproduces \cref{eq:caroli}
weighted by $\hbar\omega[n_L-n_R]$, recovering Landauer.

\subsubsection{Local heat current and conservation}

The bond-resolved heat current of Hardy~\cite{hardy_energy-flux_1963} can provide atomic insight.
Between two transport-direction slabs $i$
and $j = i + 1$,
\begin{equation}\label{eq:I_local}
  I_{i,i+1}
  = -\int_0^\infty\!\frac{d\omega}{2\pi}\,\hbar\omega\,
    \mathrm{Re}\,\mathrm{Tr}\bigl[H_{i,i+1}\,G^<_{i+1,i}(\omega)\bigr],
\end{equation}
with $H_{i,i+1}$ the inter-slab coupling block. In a converged SCBA
solution the steady-state continuity equation enforces $I_{i,i+1}$
independent of $i$ throughout the device, providing physically meaningful convergence test:
\begin{equation}
  \delta I = \max_i I_{i,i+1} - \min_i I_{i,i+1}
\end{equation}
measures the violation of energy conservation by approximations in
$\Sigma_s$ (incomplete SCBA iteration, Hilbert-transform truncation,
finite frequency grid). A convergence criteron for SCBA can be
$\delta I/\bar I \lesssim \varepsilon$~\cite{luisier_atomistic_2012,guo_quantum_2020} for
appropriate $\varepsilon$ (e.g. $10^{-3}$ in \cite{guo_quantum_2020}).

Heat conservation at the level of \cref{eq:sigma_sunset_contour}
is guaranteed by $\Phi$-derivability~\cite{baym_self-consistent_1962} only at full
self-consistency. At any finite iteration count, $\delta I/\bar I$ is
non-zero. The Hilbert-transform neglect of \cref{eq:sigma_r} introduces an additional
$\mathcal{O}(\Delta\omega/\omega)$ conservation error.

\paragraph{The conservation identity.}
The statement can be made exact at the level of \cref{eq:meir_wingreen}.
Writing the $\hbar\omega$-weighted energy current \emph{into} the device from
each terminal $a\in\{L,R,s\}$ (the two leads and the scattering ``terminal''),
\begin{equation}\label{eq:Ja}
  J_a = \int_0^\infty\!\frac{d\omega}{4\pi}\,\hbar\omega\,
  \mathrm{Tr}\bigl[\Sigma_a^<(\omega)\,G^>(\omega) - \Sigma_a^>(\omega)\,G^<(\omega)\bigr],
  \qquad \Sigma_{\rm tot} = \Sigma_L + \Sigma_R + \Sigma_s,
\end{equation}
steady-state energy balance is $J_L+J_R+J_s=0$. Summing \cref{eq:Ja} over $a$
and inserting the Keldysh form \cref{eq:keldysh_eq}
$G^{\lessgtr}=G^R\Sigma_{\rm tot}^{\lessgtr}G^A$ gives the identity
\begin{equation}\label{eq:sumrule}
  J_L+J_R+J_s = \int_0^\infty\!\frac{d\omega}{4\pi}\,\hbar\omega\,D(\omega),
  \qquad
  D(\omega) = \mathrm{Tr}\bigl[\Sigma_{\rm tot}^<\bigl(
     G^R\Sigma_{\rm tot}^>G^A - G^A\Sigma_{\rm tot}^>G^R\bigr)\bigr],
\end{equation}
where $D(\omega)\equiv\mathrm{Tr}[\Sigma_{\rm tot}^<G^>-\Sigma_{\rm tot}^>G^<]$
has been rewritten as the \emph{ordering commutator} on the right using the
cyclic property of the trace. Energy conservation is thus the pointwise
statement $D(\omega)\equiv0$, which holds through either of two independent
mechanisms. \emph{(i)~Keldysh route ($\eta\to0$).} The half rule
\cref{eq:sigma_r} satisfies $\Sigma^R-\Sigma^A=\Sigma^>-\Sigma^<$ (any
Hermitian Kramers--Kronig part of \cref{eq:sigma_R_full} cancels in the
difference), and the regulator term of \cref{eq:system}, being
$\mathcal{O}(\eta)$, drops from $[G^R]^{-1}-[G^A]^{-1}$; hence
$G^>-G^<=G^R-G^A$, the two resolvent orderings in \cref{eq:sumrule} coincide,
and $D\equiv0$ for arbitrary data. \emph{(ii)~Reciprocity route (any $\eta$).}
If every matrix is complex-symmetric, $M=M^T$ --- the reciprocity of a phonon
system whose force constants obey \cref{eq:asr} --- the transpose identity
$\mathrm{Tr}\,M=\mathrm{Tr}\,M^T$ maps one trace ordering in \cref{eq:sumrule}
into the other and $D\equiv0$ identically. At finite $\eta$ \emph{and} broken
reciprocity neither route applies and $D\neq0$.

\paragraph{Total symmetry of the vertex and $\Phi$-derivability.}
Vanishing of $J_s$ \emph{alone} is more restrictive. It requires that the
scattering self-energy be the genuine functional derivative
$\Sigma_s=\delta\Phi/\delta G$ of a Luttinger--Ward functional, evaluated with
the \emph{same} $G$ on every internal line, \emph{and} that the cubic vertex be
totally symmetric under the permutation group $S_3$,
$\Phi_{\mu\nu\lambda}=\Phi_{\sigma(\mu)\sigma(\nu)\sigma(\lambda)}$ for all
$\sigma\in S_3$. Only then do the $\hbar\omega$-weighted triads of the sunset
bubble \cref{eq:sigma_guo} telescope pairwise to give
$J_s=P_{\rm in}-P_{\rm out}=0$, the energy-current conservation of the SCBA
established in~\cite{lu_coupled_2007,wang_quantum_2008} (the explicit $S_3$
telescoping is the present derivation). This
symmetry is automatic for a vertex defined as a third derivative of the
potential (\cref{eq:H_phonon})~\cite{luisier_atomistic_2012}: a vertex that is
not totally symmetric is not the third derivative of any potential, the theory
is then not $\Phi$-derivable~\cite{baym_self-consistent_1962}, and $J_s\neq0$
even at self-consistency. At an SCBA fixed point with an exactly $S_3$ vertex
the lead currents balance, $J_L=-J_R$, at \emph{any} $\eta$; away from the
fixed point (a mixed iterate, $\Sigma_s\neq\delta\Phi/\delta G[G]$) $J_s\neq0$
and the slab imbalance \cref{eq:I_local} is finite. The imbalance is therefore
a \emph{convergence} diagnostic, not a defect of the formulation.

\paragraph{The broadening $\eta$ as a ghost reservoir.}
The residual finite-$\eta$ term in \cref{eq:sumrule} has a transparent origin.
The total anti-Hermitian broadening of the propagator,
\begin{equation}\label{eq:ghost}
  i\bigl[(G^R)^{-\dagger}-(G^R)^{-1}\bigr]
   = \Gamma_L + \Gamma_R + \Gamma_s + \Gamma_\eta,
   \qquad \Gamma_a = i(\Sigma_a^R-\Sigma_a^A),
\end{equation}
shows that the numerical regulator of \cref{eq:system} adds a damping channel
$\Gamma_\eta=\mathcal{O}(\eta)$ with \emph{no} companion fluctuation term
$\Sigma_\eta^{\lessgtr}$. A conserving (Kadanoff--Baym) theory requires every
dissipative self-energy to carry a matching $\Sigma^{\lessgtr}$ fixed by
fluctuation--dissipation balance, so $\eta$ acts as a ``ghost reservoir'' at
undefined occupation. In equilibrium, or whenever $G$ is reciprocal, its in-
and out-flows cancel and \cref{eq:sumrule} vanishes; under a temperature bias
the nonequilibrium $G^<$ is no longer complex-symmetric, the bubble inherits
the asymmetry, and the product $\eta\times\mathrm{asym}$ survives in $D(\omega)$
as a small spurious current. Consequently the slab/lead imbalance does not
reach zero at finite $\eta$ under bias but floors at the value of
\cref{eq:sumrule}; the robust convergence gates are the self-energy residual
and the pairing-exact bubble balance $P_{\rm in}=P_{\rm out}$, while transport
ratios should be extrapolated to $\eta\to0$. \Cref{app:conservation} verifies
each step of this analysis numerically and quantifies the floor on production
devices.

% ----------------------------------------------------------------------
